\documentclass[11pt]{article}

\usepackage[]{acl}          
\usepackage{latexsym}
\usepackage[T1]{fontenc}
\usepackage[utf8]{inputenc}
\usepackage{booktabs}
\usepackage{amsmath}
\usepackage{times}



\title{Classifying \& Generating Clickbait Spoilers for Text:


A Study of Extractive Question Answering and Metric Reliability}


\author{Aryan Gandhi \\
  University of Waterloo \\
  \texttt{a27gandhi@uwaterloo.ca} \\
  Student ID: 200783566}

\begin{document}
\maketitle

\begin{abstract}
A clickbait headline works by giving small bits of information while holding out
the full information so that the reader has to click to find out more about it. \textit{Clickbait
spoiling} is the task of exposing that missing information by generating a short text that satisfies the curiosity induced by a clickbait post. In this report I study and discuss the two-part version of this problem from the original competition's Clickbait
Spoiling Corpus. The first part is deciding what kind of answer a headline needs (e.g. a short
\textit{phrase}, a longer \textit{passage}, or a \textit{multi}-part texts), and
then producing the answer itself, where the predicted kind tells the generator
how to respond. Our Task 1's classifier reaches a weighted $F_1$ of
\textbf{0.759} and our second task generator reaches a METEOR
score of \textbf{0.447}. Alongside these results, I highlight two
findings that I think are more useful than the ranking. The first is that the
choice of evaluation metric quietly drives every modelling decision. In this report, I compare the two common generation metrics, BLEU and METEOR, and show that a change which
looks bad under one can be exactly the change that helps under the other. The
second is that improvements measured on a small validation set often do not carry
over to the test set, and I describe simple signals that predict which ones
will. I release all of our code and experiment logs, including the approaches
that did not work and why.
\end{abstract}


\section{Introduction}
\label{sec:intro}
Clickbait headlines exploit what psychologists call a ``curiosity gap'', meaning they
reveal just enough to make a reader want the missing details and withhold that
detail to compel the user to click on the post \citep{loewenstein1994}. \textit{Clickbait spoiling}
\citep{hagen2022} is the idea of closing that gap automatically by reading the
linked article and generating text, so the reader no longer has
to click to satisfy their curosity. Following the SemEval-2023 Task~5 formulation \citep{semeval2023task5},
I approach this as two connected sub-tasks over the same data.

\textbf{Task~1} decides in what form the spoiler should be generated. Some headlines are resolved by a short \textit{phrase} (e.g. a name or a number), some
need a full \textit{passage} (a sentence or paragraph of explanation), and some
require multiple non-consecutive pieces of text. Because this is a three class classification problem, it is scored with
weighted $F_1$, which is a standard metric for imbalanced classification problems that I explain further
in Section~\ref{sec:metrics}. \textbf{Task~2} then produces the spoiler text
itself and is scored with METEOR, a metric for comparing generated text to a
reference answer. The two tasks are connected: the type predicted in Task~1 tells
the Task~2 generator how to format it's answer.

The goal was to build a solid system for both tasks and also to
understand why particular choices helped or hurt. This paper makes three main points. It is shown that a simple, well-chosen pipeline (a text classifier for
Task~1 and a question-answering model for Task~2) is competitive on both
leaderboards. Additionally, I show that the evaluation metric itself is a design decision:
optimizing for BLEU and optimizing for METEOR can pull in opposite directions,
so picking the right target matters as much as picking the right model. Finally,
I document a number of approaches that did not work and explain the reasons. Since these were as instructive as the successes. All code, configurations, and
per-experiment logs are publicly available.%
\footnote{\texttt{github.com/TODO/mse641-clickbait-spoiling} (TODO: insert repository URL).}

\section{Background and Related Work}
\label{sec:related}
Clickbait detection, which is the problem of deciding whether a headline is clickbait or not, has a long
history \citep{potthast2016}. Spoiling, actually resolving the headline, was
introduced by \citet{hagen2022}, who released the original Clickbait Spoiling Corpus that was used. The researchers framed phrase spoilers as a \textit{question answering}
problem and passage spoilers as a \textit{retrieval} problem, and reported two
findings that shaped my design. First, question-answering models clearly beat
retrieval models at this task. Second, handling each spoiler type with its own
strategy (rather than one model for all types) improves results by several
points. SemEval-2023 Task~5 \citep{semeval2023task5} turned this into the
two-task competition that I've participated in.



\paragraph{Ideas I borrowed from prior systems.}
Two SemEval-2023 systems influenced the approach. The winning entry
\citep{billybatson2023} first reduced each article to the handful of paragraphs
most relevant to the headline, using a trained ranking model to decide relevance,
and used \textit{focal loss} \citep{lin2017} (a training objective that pays extra
attention to rare classes) to handle the type imbalance. Another strong entry
\citep{chickadams2023} reported a useful and slightly surprising result. They found that feeding
the classifier only the headline and article title, and dropping the
article body, works better than feeding the whole article. They also cleaned and preprocessed the
input text before classifying. I adopted this ``title-only'' input and found that
text cleaning is in fact the single largest improvement on Task~1
(Section~\ref{sec:results}). I also tried a simple version of their relevance
step, by keeping only the article paragraphs with the most word overlap with the
headline, but this hurt accuracy. This is likely because a plain word-overlap filter is a poor judge of which paragraph actually matters, which is presumably why the winning system used
a trained ranker instead.


\paragraph{Generating spoilers.}
More recent work by \citet{pal2024} generates spoilers with LongT5
\citep{guo2022}, a text-generation model that reads an input and \textit{writes} a
new answer rather than copying a span from it, and is built to handle long
documents. They report very high scores on passage and multi-part spoilers. As
Section~\ref{sec:results} discusses, I was not able to reproduce those numbers.

\paragraph{Models and metrics I build on.}
Our models are pre-trained language encoders: BERT \citep{devlin2019}, RoBERTa
\citep{liu2019}, and DeBERTa \citep{he2021}. These are models trained on large
amounts of text that can be adapted (``fine-tuned'') to a specific task with
comparatively little data. For Task~2 I use them in the style of SQuAD
\citep{rajpurkar2016}, a benchmark where a model is given a question and a
passage and must point to the answer inside the passage. For evaluation I use
two standard text-generation metrics, BLEU \citep{papineni2002} and METEOR
\citep{banerjee2005}, whose differences turn out to matter (see Section~\ref{sec:metric}).


\section{Approach}
\label{sec:approach}
The final system for each task is described here and the full progression
of experiments is reported in Section~\ref{sec:results}.

\subsection{Task 1: Classifying the Spoiler Type}
The classifier reads the headline together with the article title (but not the
article body, following \citet{chickadams2023}) and predicts one of the three
types. I moved from a simple model to a stronger one in stages.

Our baseline is a classical text classifier. It represents each input with TF-IDF
features, which describe a text by the words it contains, weighting words by how
distinctive they are, together with \textit{character n-grams}, which are short sequences of
letters that capture stylistic cues such as punctuation, capitalization, and
number formatting. A logistic-regression model then predicts the type (three classes total). This is a
sensible baseline because clickbait style is often visible on the surface of the
headline.


I then fine-tune a pre-trained encoder (RoBERTa, later DeBERTa) for classification. Because the three types are imbalanced, I weight
the training objective so that rarer types count more. Writing $\hat{y}_c$ for the
predicted probability of class $c$ and $N_c$ for its frequency, the objective is
the class-weighted cross-entropy
\begin{equation}
  \mathcal{L}_{\text{CE}} = -\sum_{c=1}^{3} w_c\, y_c \log \hat{y}_c,
  \qquad w_c = \frac{N}{3 N_c}.
  \label{eq:wce}
\end{equation}
I also tried focal loss, a variant that further contextualizes hard, minority-class
examples. As Section~\ref{sec:results} shows, this helped on the validation set
but hurt on the test set, likely due to overfitting. Finally, I clean the input text before classifying by
removing hyperlinks, replacing hashtags and user mentions with placeholder
tokens, stripping emojis, and expanding contractions. The final Task~1 system is
a fine-tuned RoBERTa-large encoder trained with the weighted objective of
Equation~\ref{eq:wce} on cleaned, title-only input.
\subsection{Task 2: Generating the Spoiler}
I treat spoiler generation as an \textit{extractive question-answering} problem. A
question-answering model is given a question and a passage of text and predicts
the span of the passage that answers the question. I re-purpose this directly:
the clickbait headline plays the role of the question, the linked article plays
the role of the passage, and the model learns to point at the portion of the
article that resolves the headline. This is a natural fit because the answer
almost always appears verbatim in the article, in fact, for $99.6\%$ of training examples
the ground-truth spoiler is a literal substring of the article.


I fine-tuned roberta-base-squad2, which is a RoBERTa model already trained on the
SQuAD question-answering benchmark, on the ground-truth spoiler spans. Starting
from a model that already knows how to do extractive QA matters, a plain encoder
begins with an untrained answer head and performs far worse. Because some articles can
be longer than the model's input limit, I read them with a sliding window and
keep the training checkpoint that scores best on the validation set.
\paragraph{Type-conditioned decoding.}
The type predicted in Task~1 changes how I turn the model's output into a final
answer, following the type-aware idea of \citet{hagen2022}. The strategy for each type is:
\begin{itemize}
  \item \textit{phrase}: return the single best-scoring span, capped at a short length.
  \item \textit{passage}: return the \textit{entire paragraph} that contains the
  best span, rather than the span alone.
  \item \textit{multi}: return the several best non-overlapping spans, joined together.
\end{itemize}

The passage rule comes from how METEOR scores answers. Because METEOR rewards
when the output includes the reference text (recall) more than being concise (precision), returning
the whole paragraph, which is likely to contain the reference text, beats returning a
short fragment of it, even though the paragraph includes extra words. I found
that for passage spoilers the score is almost entirely determined by whether I
pick the right paragraph. If $a$ is the fraction of the time I select a paragraph
that actually contains the answer, then
\begin{equation}
  \text{METEOR}_{\text{passage}} \approx 0.77\,a + 0.11\,(1-a),
  \label{eq:passage}
\end{equation}
that is, a correctly chosen paragraph scores about $0.77$ and a wrong one about
$0.11$. Section~\ref{sec:error} returns to this.

\section{Experiments and Results}
\label{sec:results}

\subsection{Dataset}
I used the Clickbait Spoiling Corpus as provided
by the competition, with a fixed split of $3{,}200$ training, $400$ validation,
and $400$ test examples (the test labels and spoilers are withheld and scored only through the
leaderboard). The types are imbalanced in training set with the class distribution being $42.7\%$ \textit{phrase}, $39.8\%$
\textit{passage}, and $17.5\%$ \textit{multi}. Each example contains the
clickbait post, the linked article, and for the training and validation splits,
the ground-truth spoiler and its type are included.

\subsection{Evaluation Metrics}
\label{sec:metrics}
The two tasks use different metrics because they are different kinds of problems.
Task~1 is classification, so I score it with \textbf{weighted $F_1$}. The $F_1$
score is the harmonic mean of precision and recall for a class, however, the weighted
version averages the per-class $F_1$ scores in proportion to how often each type
occurs. For this problem weighted $F_1$ score should be used rather than plain accuracy because the types are imbalanced and because of this the
accuracy can look high while the model quietly fails on the rare
classes (e.g. multi), whereas weighted $F_1$ only looks good if the model handles all three
types reasonably well despite this.

Task~2 is text generation, so we need a metric that compares a produced answer to
a reference text (i.e. ground truth). I use METEOR, the competition's official metric. Unlike
exact match, METEOR gives partial credit, it matches words that share a stem or
are synonyms and it weights recall (how much of the reference the answer covers)
above precision (how concise the answer is while matching the reference). This suits spoiling where the point
is to convey the withheld information, not to reproduce the exact wording. I also
compute BLEU, an older overlap-based metric that focuses on precision and use it
as a point of comparison. Section~\ref{sec:metric} explains why the comparison is
informative.

\subsection{Experimental Setup}
All experiments use a fixed random seed and the same evaluation harness, so that
every run is logged with its configuration and score. Transformers are fine-tuned
with the AdamW optimizer, a linearly decaying learning rate, and gradient
clipping. The Task~1 classifiers use an input length of $256$ tokens, a learning
rate of $1\mathrm{e}{-5}$, and up to five or six epochs. The Task~2 model uses a
learning rate of $3\mathrm{e}{-5}$ and three epochs. For every model I keep the
checkpoint that scores best on the validation set. Training was done on a personal
laptop, which limited the largest model I could use to RoBERTa-large.
Section~\ref{sec:conclusion} reflects on this.

\subsection{Task 1 Results}
Table~\ref{tab:task1} traces the progression. The largest early gain comes from
adding character n-grams to the baseline ($+0.078$), which confirms that a lot of
the signal is stylistic and lives on the surface of the headline. Moving to a
fine-tuned transformer gives far better results and our best result of
\textbf{0.759} on the test set comes from RoBERTa-large trained with
the weighted objective on cleaned input.

\begin{table}[t]
\centering\small
\begin{tabular}{@{}llcc@{}}
\toprule
\# & System & Val $F_1$ & Test $F_1$ \\
\midrule
001 & TF-IDF + logistic regression & 0.527 & 0.540 \\
003 & \ + headline character n-grams & 0.626 & -- \\
007 & \ + RoBERTa (blended) & 0.695 & 0.706 \\
012 & DeBERTa-v3-base & 0.700 & 0.729 \\
013 & RoBERTa-large + focal loss & 0.716 & 0.703 \\
014 & RoBERTa-large + weighted loss & 0.742 & 0.732 \\
\textbf{015} & \ \textbf{+ text cleaning} & \textbf{0.749} & \textbf{0.759} \\
\bottomrule
\end{tabular}
\caption{Task~1 progression (weighted $F_1$). Our best test score is 0.759 and a baseline of 0.540.}
\label{tab:task1}
\end{table}

Two details are worth noting. Blending the linear model with the transformer
helped RoBERTa slightly but \textit{hurt} the stronger DeBERTa model, because by that
point the transformer already captured the stylistic cues the linear model was
contributing. And the \textit{multi} class was the hardest throughout (its
per-class $F_1$ is $0.68$, against $0.77$ for the other two types), since it is
the rarest and is easily confused with the passage type.

\subsection{Task 2 Results}
Table~\ref{tab:task2} traces Task~2. The two biggest gains are fine-tuning the
question-answering model on the ground-truth spoiler spans and adding
type-conditioned decoding. Returning the full paragraph for passage spoilers and
tuning the per-type decoding reached a test METEOR score of \textbf{0.447}.

\begin{table}[t]
\centering\small
\begin{tabular}{@{}llcc@{}}
\toprule
\# & System & Val & Test \\
\midrule
002 & Question answering, no fine-tuning & -- & 0.214 \\
005 & Fine-tuned QA + type decoding & -- & 0.440 \\
008 & \ + full-paragraph passages & 0.426 & 0.445 \\
\textbf{010} & \ \textbf{+ tuned per-type decoding} & \textbf{0.471} & \textbf{0.447} \\
\bottomrule
\end{tabular}
\caption{Task~2 progression (METEOR). Our best test score is 0.447 with a baseline of 0.214.}
\label{tab:task2}
\end{table}

\subsection{The Metric Matters: BLEU versus METEOR}
\label{sec:metric}
The competition scores METEOR, but BLEU is another popular and familiar NLP metric, I
tracked both for comparison which turned out to be very revealing. When I changed the
passage decoder to return the whole paragraph (experiment 008), METEOR
\textit{rose} by $0.018$ while BLEU \textit{fell} from $28.5$ to $22.2$. The reason is
exactly the trade-off between the two metrics, since returning the whole paragraph adds
words, which BLEU penalizes as imprecise, but which METEOR rewards for covering
more of the reference. Had I been optimizing for BLEU, I would have rejected
the change that most improved the metric the competition was actually using for grading. The lesson
is that the metric is not a neutral scorer, it silently decides which changes count as progress and is opinionated in what matters. Additionally, an argument can be made that precision should be more equally weighted to recall than METEOR currently allows, to encourage more concise spoilers.

\subsection{Validation Gains Do Not Always Transfer}
\label{sec:transfer}
A recurring theme was that an improvement on the $400$-example validation set did
not reliably carry over to the test set, and occasionally reversed
(Table~\ref{tab:transfer}).

\begin{table}[t]
\centering\footnotesize
\setlength{\tabcolsep}{1pt}
\begin{tabular}{@{}lccl@{}}
\toprule
Change & $\Delta$Val & $\Delta$Test & Outcome \\
\midrule
Task~2 decoding sweep & $+0.044$ & $+0.003$ & barely transferred \\
Task~1 focal loss & $+0.015$ & $-0.027$ & reversed \\
Task~1 weighted loss & $+0.041$ & $+0.003$ & small \\
Task~1 text cleaning & $+0.007$ & $\mathbf{+0.027}$ & test beat val \\
\bottomrule
\end{tabular}
\caption{Validation improvements transferred inconsistently to the test set.}
\label{tab:transfer}
\end{table}

Two patterns explain most of this. First, a balanced set of predictions was a
better sign of real improvement than the validation score itself. Focal loss
raised validation $F_1$ by predicting the rare \textit{multi} type more
aggressively, but on the test set those extra guesses were mostly wrong and the
score fell; the plainer weighted-loss model, whose predictions were more balanced,
transferred cleanly. Second, changes that made the input cleaner or the method
more robust, rather than fitting the small validation set, transferred best, since text
cleaning was the only change whose test gain actually exceeded its validation gain.

\subsection{Error Analysis}
\label{sec:error}
Because passages are where most of our lost points sit, I looked at them closely
on the $154$ passage-type validation examples, and I found a lot about where the system stands and what would move it.

\paragraph{The passage bottleneck: picking the right paragraph.}
Two facts stand out. The ground-truth spoiler is present somewhere in the article
every single time, so nothing is actually missing from the input. Yet the
paragraph our system returns contains that spoiler only $54.5\%$ of the time. To put it
plainly, the answer is always there and we simply reach for the wrong paragraph
almost half the time. This is also where Equation~\ref{eq:passage} comes from. To explain it further, I measured the average METEOR separately for the two
cases and found that when the returned paragraph does contain the spoiler the
score averages $0.77$, and when it does not it averages only $0.11$. The overall
passage score is just the mixture of those two outcomes weighted by how often each
happens, which is exactly Equation~\ref{eq:passage}. At our $54.5\%$ selection
rate it predicts about $0.47$, which matches the $0.46$ we actually observe.

\paragraph{Two separate gaps to a perfect score.}
That decomposition shows the distance to a perfect $1.0$ is really two different
gaps. The large one is \textit{selection}, which involves raising the fraction of correctly
chosen paragraphs from $0.55$ toward $1.0$, this is worth the most by far. The smaller
one is \textit{granularity}, even when we pick the right paragraph, we return the
\textit{whole} paragraph, which is much longer than the spoiler itself, only about
a third of the paragraph's length on average. METEOR rewards us for covering the
reference but punishes us for the surrounding sentences we drag along, which is why a
correctly chosen paragraph averages $0.77$ rather than $1.0$. So even flawless
selection would cap this strategy near $0.77$. The true ceiling is slightly lower,
about $0.74$, because in roughly $4.5\%$ of cases no single paragraph contains the
entire spoiler at all. This is because the answer sometimes can span a paragraph break or spills over from
the article title, and a one-paragraph answer can never capture all of it.
Trimming the paragraph down to just the spoiler sentence would recover some
precision but is a tough task. It also risks cutting part of the answer and since METEOR punishes lost
coverage more heavily than extra words, over-trimming tends to score worse
than leaving the paragraph whole. That reward system is why we return the full paragraph.

\paragraph{Why the selection goes wrong.}
The reason we miss the right paragraph so often traces back to how the model was
trained. A question-answering model learns to find short, precise answer spans,
which is ideal for phrase spoilers but is the wrong instinct for passages. Because it is asked
to locate a whole passage, it anchors on the single most answer-like sentence and
when a plausible but incorrect sentence appears earlier in the text, it commits
to that paragraph instead. The four selection methods in
Section~\ref{sec:negative} were all attempts to override this instinct with a
separate relevance signal, and since none of them beat simply trusting the model's
own span, the more promising fix looks like a stronger extractive model rather
than a ranker one.

\paragraph{How classifier errors flow into Task 2.}
Because the Task~1 prediction determines the Task~2 decoding rule, a misclassified
type applies the wrong strategy to an input text. To measure how much this costs, I
re-ran Task~2 using the true types instead of our $69.5\%$-accurate
predictions, and validation METEOR rose only from $0.471$ to $0.510$. That
$0.039$ is the most a perfect classifier could ever add and a realistic
classifier would recover only part of it, so good type prediction does add value but only slightly. The errors that do occur are often concentrated on the multi type,
which is both the rarest and the easiest to confuse with a passage, which is aligned with it's
status as the hardest class in Task~1.

\paragraph{Two failure modes the metric exposes.}
Some of our lowest scores say more about the metric than the model. Phrase
spoilers show this most clearly because the answer is often a number or a name and if
our wording differs at all from the reference we can score zero even when the
meaning is right. Multi spoilers show a different story, where the real difficulty
is deciding \textit{how many} pieces to return and \textit{which} ones. For example, recovering
two of three correct snippets earns partial credit, but the model has no reliable
signal for the true count and tends to guess.

\paragraph{A few concrete examples.}
These quirks are easier to see in real outputs. For a \textit{phrase} spoiler, the
post ``You won't believe what Samir Nasri has done\ldots'' has the ground-truth
spoiler \textit{dying his hair sky blue}, which the model recovers exactly. A
revealing phrase failure shows the metric strictness above: for a post
about how much to tip, the ground truth is \textit{20\%} but the model answers
\textit{between \$5 and \$15}, which is the right idea yet shares no words with the
reference and so scores zero. A \textit{passage} example shows the selection
problem directly: the ground-truth spoiler is \textit{some of the plot elements
are so disturbing that they are making him feel sick}, but the model returns a
different, shorter piece of the article (\textit{too dark}) and misses it.
Finally, for a \textit{multi} spoiler whose answer has three parts (\textit{Alan
Rickman \& Rupert Grint CBGB}), the model recovers two of them, which is a partial success.

\subsection{What Did Not Work}
\label{sec:negative}
Several reasonable-looking ideas failed and I think the reasons are as useful as
the successes (Table~\ref{tab:negative}). The clearest pattern concerns paragraph
selection for passages. I tried four different ways to pick the right paragraph
(word-overlap similarity, aggregating the model's span scores per paragraph, a
combination of the two, and a separately trained ranking model), and all four did
worse than simply trusting the question-answering model's own best span. This
independently reproduces \citet{hagen2022}'s finding that question answering beats
retrieval for this task. I also tried generating spoilers with a text-generation
model instead of extracting them, but a smaller model reached only $0.31$
METEOR on passage and multi spoilers, against $0.42$ for extraction. 

\begin{table}[t]
\centering\small
\begin{tabular}{@{}p{3.3cm}p{3.7cm}@{}}
\toprule
What I tried & Why it did not help \\
\midrule

Keeping only word-overlap paragraphs (Task~1) & word overlap is a poor judge of relevance \\\\

Focal loss (Task~1) & over-predicted the rare class; helped val, hurt test \\\\

Generating instead of extracting (Task~2) & tended to paraphrase; extraction scored higher \\\\

Four paragraph-selection methods (Task~2) & all lost to the QA model's own span \\\\

Returning several paragraphs (Task~2) & extra coverage was cancelled by extra noise \\\\

\bottomrule
\end{tabular}
\caption{Approaches that did not improve results, and the reason in each case.}
\label{tab:negative}
\end{table}

\section{Conclusion and Future Work}
\label{sec:conclusion}
I built a two-task clickbait-spoiling system, a text classifier for the spoiler
type and a fine-tuned question-answering model for the spoiler text, that performed well across both tasks in the competition. I came away with two lessons that generalize beyond this competition. The evaluation metric is
arguably the most important design choice, since optimizing for BLEU and optimizing for METEOR reward
different behaviour, so choosing the right target is as consequential as choosing
the right model. And a small validation set is an unreliable guide, since improvements
that come from cleaner inputs or more balanced predictions transfer to the test
set far more dependably than improvements squeezed out by tuning to the validation
numbers. However, this result could be different if the validation set was very large.

The clearest path forward on Task~2 is better paragraph selection for passage
spoilers, since perfect selection would raise our score from about $0.46$ to
$0.74$. However, keep in the mind every selection method I tried lost to the extractive model, so
a stronger extractive model is the more promising direction than a
ranker. A larger classifier such as DeBERTa-v3-large, combined with our text
cleaning, is a natural next step for Task~1. Additionally, in future work I would have trained
these larger models on more powerful cloud GPUs from the start. I did not realize early enough how limiting and slower training is on a personal laptop than on the cloud infrastructure used in prior work.


\bibliography{references}

\end{document}
